\begin{textsample}{POS Dim 2 – df – Score 18.00 – df/df\_em\_54.txt}  \label{ex:f2_pos_129}
\textbf{CURSOS} DE EDUCAÇÃO \textbf{PROFISSIONAL} E \textbf{TÉCNICA} DE NÍVEL MÉDIO - \textbf{Cargas} \textbf{horárias} de 800 ( oitocentas ), 1.000 ( mil ) ou 1.200 ( mil e duzentas ) horas ( CNCT ). - \textbf{Habilitação} \textbf{profissional} para a vida e para o mundo do trabalho, desenvolvendo as capacidades para utilizar, desenvolver ou adaptar tecnologias, compreendendo as implicações delas decorrentes, bem como suas relações com o processo produtivo e com a sociedade. - Diploma com validade nacional para \textbf{cursos} devidamente registrados no SISTEC. - Os \textbf{cursos} \textbf{técnicos} de nível médio poderão ser \textbf{ofertados} de forma subsequente, para quem já concluiu o Ensino Médio, ou de forma integrada, na mesma instituição ; ou concomitante, na mesma, ou em distintas instituições. \textbf{CURSOS} DE \textbf{QUALIFICAÇÃO} \textbf{PROFISSIONAL} OU DE FORMAÇÃO INICIAL E CONTINUADA ( FIC ) - \textbf{Carga} \textbf{horária} mínima de 160 ( cento e \textbf{sessenta} ) horas. - \textbf{Qualificação} \textbf{profissional} para a vida e para a inserção ou reinserção no mundo do trabalho. - Reconhecidos por meio de certificado.
PRINCÍPIOS DA EDUCAÇÃO \textbf{PROFISSIONAL} E TECNOLÓGICA ( EPT ) PRINCÍPIOS \textbf{ORIENTADORES} Os princípios que orientam a EPT em todas as suas fases de organização, formas e modalidades de \textbf{oferta}, no âmbito do Sistema de Ensino do Distrito Federal, com base na Resolução nº 06/2012 – CNE ( BRASIL, 2012 ), são : - formação integral do estudante, garantindo articulação entre a formação no Ensino Médio e a formação \textbf{técnica} específica ; - pleno desenvolvimento para a vida social e \textbf{profissional}, com formação voltada para o \textbf{fortalecimento} de valores cidadãos ; - respeito aos conhecimentos prévios dos sujeitos da aprendizagem, promovendo a indissociabilidade entre educação e prática social ; - indissociabilidade entre teoria e prática no processo de ensino e aprendizagem ; - envolvimento das múltiplas dimensões do eixo tecnológico do \textbf{curso} e das ciências e tecnologias a ele vinculadas ; - articulação com os arranjos produtivos locais, com as demandas da comunidade escolar e com o desenvolvimento socioeconômico-ambiental dos territórios ; - reconhecimento da diversidade das pessoas com deficiência, transtornos globais do desenvolvimento e altas habilidades, em \textbf{regime} de acolhimento ou internação e em \textbf{regime} de privação de liberdade ; - reconhecimento das identidades étnico-raciais dos povos indígenas, dos quilombolas e das identidades de gênero, bem como da diversidade de orientação sexual, das populações do campo e das pessoas em situação de rua ; - autonomia da Instituição Educacional na construção e revisão de seu Projeto Político Pedagógico ( PPP ), desde que respeitados o processo participativo, a \textbf{legislação} e as normas educacionais ; - descrição de \textbf{perfis} \textbf{profissionais} de conclusão de \textbf{curso} que contemplem conhecimentos, competências e saberes \textbf{profissionais} requeridos pela natureza do trabalho, pelo desenvolvimento tecnológico e pelas demandas sociais, econômicas e ambientais ; - colaboração entre os entes federados, incluindo parcerias com as \textbf{entidades} da sociedade civil, visando à melhoria dos indicadores educacionais.

% matched lemmas: carga, curso, entidade, fortalecimento, habilitação, horário, legislação, oferta, ofertar, orientador, perfil, profissional, qualificação, regime, sessenta, técnico
\end{textsample}
